\section{Introduction}\label{sec:intro}

Machine learning is now a core empirical methodology in quantitative finance
\citep{GuKellyXiu2020}, which makes its results unusually dependent on code: the
method is the implementation, and a description in prose rarely determines it.
A reader who wants to build on such a result therefore has a concrete task:
obtain the authors' code, install what it needs, run it, and see whether the
numbers come out. Almost everything the reproducibility literature measures is
upstream of that task. Instruments count whether a paper states that code is
available, whether a link is present, whether data are described, whether an
environment is specified. These are properties of the manuscript. They are cheap
to score at scale, they are useful, and they are not the same thing as being able
to run the code.

The gap matters because the two can diverge in both directions. A paper can
disclose exemplary metadata, a repository link, a license, a
\texttt{requirements.txt}, and still be unusable, because the pinned dependencies
no longer resolve, because the entry point expects a data file that was never
shipped, or because the repository was deleted after publication. Conversely a
sparse repository with no formal environment specification may run without
incident.

Finance does have a direct measurement of computational reproducibility.
\citet{Perignon2024} ran the code of 168 research teams on the same raw data and
found that it regenerates exactly the same results 52\% of the time. That figure
is a best case by construction: the code was written for a coordinated study, the
data were supplied, and the teams were reachable. A reader of a published paper
has none of those advantages. They have a URL.

This paper measures the cold-start case. We take a corpus of machine-learning papers in
quantitative finance, extract every code URL the authors advertise, and attempt
to run what we find: resolve the link, clone the repository, build the declared
environment inside a container, locate an entry point, and execute it under a
fixed time budget. What we record for each artifact is the furthest point it
reached and the reason it stopped.

\subsection{The instrument}\label{sec:intro_instrument}

Our first contribution is the measurement design itself. Ordered scales for
artifact quality are not new, and \Cref{sec:novelty} sets out the ones that
precede us; what they have in common is that a person assigns the level. A score
assigned by a reader also conflates failures with nothing in common when it is
reduced to an aggregate: a deleted repository and an unresolvable dependency both
cost a point, and the total cannot distinguish a hosting problem from a packaging
problem. We use an ordered, cumulative ladder of eleven rungs, decided by a
program rather than a reader,
(\Cref{sec:ladder}), running from ``the paper makes no code-availability
claim'' to ``an emitted number matches the paper's headline result''. Each
artifact is pushed as far as it will go; the distribution of stopping points is
the attrition curve, and the reasons for stopping form a second, orthogonal
table.

Two properties of the design carry most of its weight. Because the rungs are
cumulative, the curve is monotone by construction, and a violation of
monotonicity is a detectable scoring inconsistency rather than a quirk of the
data; the audit script checks it on every build. Because failure reasons are
drawn from a closed vocabulary rather than free text, the failure mix at each
rung is itself reportable, so we can say not only how many artifacts stopped at
environment construction but what fraction of those were caused by a package
disappearing from the index rather than by a version pin that has aged out.

\subsection{Measuring the measurement}\label{sec:intro_reflexive}

A study of this kind carries an unusual risk: a defect in the harness produces a
number that looks like a finding. The failure modes are not hypothetical. During
development, a container backend that started a fresh container for the install
step and another for the execution step produced a clean, plausible result in
which no artifact in the corpus executed at all; every dependency installed in
the first container had been discarded before the second one ran. The number was
publishable-looking and entirely an artifact of our own code. It was caught
because the failures took a tenth of a second, which is not how real
machine-learning code fails.

We therefore treat the harness as an object of measurement too. Failures
attributable to our tooling carry their own reason code and are excluded from
every claim about the literature. Conditions that indicate the harness is
misconfigured, rather than that a paper is deficient, raise an error instead of
being recorded as data. The entry-point heuristic, the most discretionary part of
the pipeline, is run in two configurations and both curves are reported, so a
reader can see how much of the result is the literature and how much is our own
conservatism. \Cref{sec:discussion} returns to what this does and does not
license us to conclude.

\subsection{Contributions}\label{sec:contributions}

\begin{itemize}[leftmargin=*]
    \item The Reproducibility Attrition Ladder: an ordered, cumulative scale for
          execution-based reproducibility, with a closed failure vocabulary that
          makes the reasons for attrition reportable rather than anecdotal.
    \item An execution study of a corpus of arXiv quantitative-finance
          machine-learning papers, reporting the attrition curve, the failure mix
          at each rung, and a sensitivity analysis over the one heuristic that
          materially affects the upper rungs.
    \item An open, installable harness that performs the measurement, so the
          study can be repeated, extended to other corpora, and re-run as the
          literature changes. The per-paper outcome records are released in full.
\end{itemize}

\subsection{Roadmap}\label{sec:roadmap}

\Cref{sec:related} places the work against the reproducibility literature in
science generally, in empirical finance, and in machine learning.
\Cref{sec:ladder} defines the ladder and the execution environment.
\Cref{sec:results} reports the corpus, the attrition curve, and the failure mix.
\Cref{sec:discussion} discusses what the curve implies for readers, authors, and
editors, and is explicit about the limits of the claim.
